\begin{acknowledgement}

行文至此，意味着四年的大学时光即将画上句号。回首来路，有代码报错时的焦虑，有仿真通过时的欣喜，更有在时代浪潮下寻找个人定位的迷茫与求索。这篇论文的完成，不仅是对我本科四年专业知识的一次大考，更是我走向独立思考、拥抱更广阔技术世界的一次洗礼。在此，我要向所有在这段旅程中给予我帮助与启迪的人致以最诚挚的谢意。

首先向我的指导老师表达最深切的感激。熊老师在论文的选题、资料搜集及不同模块的搭建过程中，都给予了我重要的支持。此外，熊老师还为我未来的发展方向提供了非常实际的帮助。老师多方面的指导是我完成论文、不断努力的后盾。

其次，感谢大学四年里遇到的所有关心支持我的人。各科教授在课堂上的传道授业，为我打下了扎实的电气工程理论基础；辅导员在生活中的关心与引导，让我在面对升学压力与职业选择时，能够保持清醒与坚定；舍友们在面对困难时的互帮互助，让我在学习与生活中有了朋友与依靠。同时，我也要感谢开源社区中那些未曾谋面的开发者与研究者们。在这四年求学之旅中，是你们无私分享的知识、讨论和代码片段，让我在无数个深夜的调试中看到了方向。大家的帮助将永远激励我在未来的工程实践中不断前行。

此外，我还要特别感谢我的父母。无论我做出何种人生选择，遇到了何种困难，父母始终是我最坚定的支持者。

最后，我想感谢那个在 Simulink 的建模前反复调参、在 MATLAB 脚本里逐行修改代码的自己。感谢自己在面对复杂模型时的没有放弃，感谢自己在面对未知技术领域的勇于探索。时代的风云变幻或许非个人所能掌控，但敲击键盘、构建逻辑的能力却永远属于自己。愿我们在未来的岁月里，都能保持对技术的热爱、对世界的好奇，以及在喧嚣中坚守本心的勇气。带着这份沉甸甸的感激，我将整理行装，勇敢地迈向人生的下一段旅程。
\end{acknowledgement}